\begin{table}[H]
\centering
\caption{\label{tab:open-ended-frame-effects}Framing differences in coded post-choice explanations}
\centering
\fontsize{8}{10}\selectfont
\begin{threeparttable}
\begin{tabular}[t]{lrrrrrr}
\toprule
Category & N & Gain \% & Loss \% & Loss - gain [95\% CI] & p & Holm p\\
\midrule
No stated reason & 1752 & 8.0 & 6.9 & -1.1 [-3.5, 1.4] & 0.390 & 1.000\\
Equal expected value & 1752 & 7.0 & 6.6 & -0.4 [-2.8, 2.0] & 0.738 & 1.000\\
Certainty or risk aversion & 1752 & 23.7 & 26.7 & 2.8 [-1.3, 6.9] & 0.177 & 1.000\\
Risk offers chance all survive & 1752 & 33.3 & 31.1 & -2.0 [-6.4, 2.3] & 0.361 & 1.000\\
Saves more people & 1752 & 18.6 & 19.4 & 0.9 [-2.8, 4.6] & 0.639 & 1.000\\
Easier to understand & 1752 & 0.2 & 0.8 & 0.6 [-0.1, 1.3] & 0.086 & 0.684\\
Other & 1752 & 6.0 & 5.5 & -0.5 [-2.7, 1.7] & 0.665 & 1.000\\
Better choice without a reason & 1752 & 3.1 & 2.9 & -0.3 [-1.9, 1.3] & 0.733 & 1.000\\
\bottomrule
\end{tabular}
\begin{tablenotes}[para]
\item Note: Percentages use Anne's primary code among attention-check passers with a valid instrument and framing outcome. Models adjust for wording pair and use HC2 standard errors. Codes classify explanations after the choice; they are descriptive and do not identify a causal mechanism.
\end{tablenotes}
\end{threeparttable}
\end{table}
